\chapter{Cơ sở lý thuyết và nền tảng toán học}
\label{chap:co_so_ly_thuyet}

Chương này trình bày hệ thống nền tảng toán học, giải tích đồ thị phổ và lý thuyết học sâu làm cơ sở khoa học cho kiến trúc Two-Stage Pipeline của đề tài. Cấu trúc chương được tổ chức tập trung vào hai trụ cột công nghệ cốt lõi:
\begin{enumerate}
    \item \textbf{Mạng xương sống EfficientNet và Mô hình hồi quy đếm xe trực tiếp (Giai đoạn 1):} Trình bày cơ sở toán học của bài toán đếm xe qua ánh xạ trực tiếp từ không gian ảnh tĩnh sang vector mật độ phương tiện theo phân lớp; phân tích nguyên lý mở rộng đồng thời (Compound Scaling), cấu trúc khối tích chập MBConv tích hợp cơ chế chú ý kênh Squeeze-and-Excitation, tầng nén đặc trưng Bottleneck MLP và hàm mất mát bền vững Huber Loss.
    \item \textbf{Mạng nơ-ron đồ thị không -- thời gian nhận biết thời gian TA-STGCN (Giai đoạn 2):} Trình bày mô hình hóa mạng lưới giao thông trên đồ thị phi Euclid; phân rã giải tích đồ thị phổ với toán tử xấp xỉ Chebyshev, cơ chế tích chập thời gian có cổng GLU; và kiến trúc đề xuất TA-STGCN tích hợp cơ chế Self-Attention theo thời gian đa đầu (Multi-Head Temporal Self-Attention) giải quyết hiện tượng lan truyền sai số cảm biến.
\end{enumerate}

% ==============================================================================
\section{Mạng xương sống EfficientNet và Mô hình hồi quy đếm xe trực tiếp}
\label{sec:kien_truc_efficientnet}

\subsection{Tổng quan nguyên lý và bài toán đếm xe qua hồi quy trực tiếp}
\label{sec:nguyen_ly_hoi_quy_dem_xe}

Trong giám sát giao thông đô thị bằng thị giác máy tính, khung làm việc truyền thống thường dựa trên mô hình \textit{Tracking-by-Detection} (kết hợp các mạng phát hiện hộp bao như Faster R-CNN, YOLO với thuật toán bám vết liên khung hình như SORT, DeepSORT). Tuy nhiên, tại các đô thị đặc thù ở Việt Nam với tỷ trọng phương tiện cá nhân hai bánh chiếm trên $80\%$, mô hình này bộc lộ những giới hạn căn bản:
\begin{enumerate}
    \item \textbf{Che khuất tương hỗ dày đặc (Severe Mutual Occlusion):} Dòng xe máy di chuyển thành cụm với cự ly rất hẹp, dẫn đến hiện tượng phần lớn thân xe bị che khuất. Các thuật toán sinh vùng bao thường xuyên bỏ sót phương tiện hoặc gộp nhiều xe máy liền kề vào cùng một hộp định vị.
    \item \textbf{Nút thắt tính toán của thuật toán triệt tiêu cực đại cục bộ (NMS):} Khi mật độ xe trong khung hình vượt quá ngưỡng cao (trên 100 phương tiện), số lượng hộp bao dự đoán sinh ra tăng đột biến. Thuật toán NMS với độ phức tạp tính toán tiệm cận $\mathcal{O}(M^2)$ (với $M$ là số hộp bao ứng viên) tạo thành điểm nghẽn trễ lớn, không đáp ứng yêu cầu xử lý thời gian thực trên các thiết bị biên.
\end{enumerate}

Để khắc phục triệt để các hạn chế trên, đề tài xây dựng bài toán đếm xe theo hướng \textbf{Hồi quy trực tiếp từ ảnh tĩnh} (Direct Density Regression). Về mặt toán học, bài toán được phát biểu dưới dạng tìm ánh xạ phi tuyến tham số hóa $f_\theta$:
\begin{equation}
    f_\theta: \mathbf{I} \in \mathbb{R}^{3 \times H \times W} \longrightarrow \hat{\mathbf{c}} = \begin{bmatrix} \hat{c}_{\text{car}} \\ \hat{c}_{\text{bike}} \end{bmatrix} \in \mathbb{R}^2
\end{equation}
trong đó $\mathbf{I}$ là tensor biểu diễn ảnh chụp tĩnh từ camera quan sát có kích thước $H \times W$ trên 3 kênh màu RGB; $\hat{\mathbf{c}}$ là vector số thực hai chiều ước lượng số lượng ô tô và xe máy hiện diện trong khung cảnh.

Mô hình học trực tiếp mối quan hệ giữa các mẫu hoa văn mật độ điểm ảnh toàn cục và tổng lượng phương tiện mà không cần định vị tọa độ từng xe riêng lẻ. Cách tiếp cận này triệt tiêu hoàn toàn khâu hậu xử lý NMS, đưa độ trễ suy luận về mức cố định (dưới 35 ms trên GPU biên), đồng thời duy trì độ chính xác cao ngay cả khi các phương tiện bị che khuất nghiêm trọng.

\subsection{Nguyên lý mở rộng đồng thời (Compound Scaling) của EfficientNet}

Để đảm nhận vai trò mạng xương sống (backbone) trích xuất đặc trưng hình ảnh cho bài toán hồi quy đếm xe, mô hình đòi hỏi khả năng biểu diễn phong phú song phải tối ưu hóa tài nguyên tính toán. Mạng nơ-ron tích chập truyền thống thường cải thiện độ chính xác bằng cách mở rộng độc lập một trong ba chiều: chiều sâu mạng ($d$), chiều rộng số kênh ($w$), hoặc độ phân giải ảnh đầu vào ($r$). Tuy nhiên, Tan và Le \cite{efficientnet} đã chỉ ra rằng việc tăng đơn lẻ từng chiều nhanh chóng dẫn đến hiện tượng suy giảm hiệu suất biên (diminishing returns).

Nhằm khắc phục nhược điểm trên, nguyên lý \textbf{Compound Scaling} điều chỉnh đồng bộ cả ba chiều thông qua một hệ số mở rộng duy nhất $\phi$:
\begin{equation}
    d = \alpha^\phi, \quad w = \beta^\phi, \quad r = \gamma^\phi
\end{equation}
với ràng buộc cân bằng mức tăng trưởng phép tính:
\begin{equation}
    \alpha \cdot \beta^2 \cdot \gamma^2 \approx 2, \quad \alpha \ge 1, \ \beta \ge 1, \ \gamma \ge 1
\end{equation}
trong đó $\alpha, \beta, \gamma$ là các hằng số phân bổ xác định qua tìm kiếm kiến trúc trên mạng cơ sở EfficientNet-B0. Bằng cách mở rộng cân đối theo $\phi$, kiến trúc duy trì tỷ lệ tối ưu giữa kích thước trường thụ cảm và dung lượng kênh, đạt độ chính xác trích xuất cao trên cùng mức chi phí tính toán (FLOPs).

\subsection{Cấu trúc khối MBConv và cơ chế chú ý kênh Squeeze-and-Excitation}

Đơn vị cấu trúc cốt lõi của EfficientNet là khối \textbf{MBConv (Mobile Inverted Bottleneck Convolution)}, tích hợp ba kỹ thuật tối ưu:

\textbf{1. Tích chập tách rời theo chiều sâu (Depthwise Separable Convolution):}
Phép toán phân tách tích chập tiêu chuẩn thành hai bước độc lập: \textit{Depthwise Convolution} (áp dụng bộ lọc $k \times k$ riêng cho từng kênh) và \textit{Pointwise Convolution} (tích chập $1 \times 1$ để kết hợp các kênh). Tỷ số chi phí tính toán đạt được là:
\begin{equation}
    \frac{\text{FLOPs}_{\text{separable}}}{\text{FLOPs}_{\text{standard}}} = \frac{1}{C_{\text{out}}} + \frac{1}{k^2}
\end{equation}
Với kernel kích thước $k = 3$, phương pháp này cắt giảm khoảng $85\% - 89\%$ lượng phép tính so với tích chập 2D thông thường mà vẫn bảo toàn năng lực biểu diễn không gian.

\textbf{2. Khối phần dư nghịch đảo (Inverted Residuals):}
Khác với khối Bottleneck truyền thống (nén kênh ở giữa), MBConv mở rộng số kênh ở tầng giữa lên gấp $t$ lần ($t=6$) trước khi áp dụng Depthwise Convolution, sau đó chiếu ngược về số kênh thấp kết hợp đường dẫn tắt (shortcut connection), giúp bảo toàn tối đa thông tin ngữ nghĩa.

\textbf{3. Cơ chế chú ý kênh Squeeze-and-Excitation (SE):}
Khối SE tự động tái hiệu chuẩn trọng số các kênh đặc trưng qua hai pha:
\begin{align}
    \mathbf{z} &= \text{GAP}(\mathbf{X}) \in \mathbb{R}^C \quad &\text{(Pha Squeeze)} \\
    \mathbf{s} &= \sigma\left( \mathbf{W}_2 \cdot \text{SiLU}(\mathbf{W}_1 \mathbf{z}) \right) \in \mathbb{R}^C \quad &\text{(Pha Excitation)}
\end{align}
với $\mathbf{W}_1 \in \mathbb{R}^{\frac{C}{r} \times C}$ (tỷ lệ co thắt $r=4$), $\mathbf{W}_2 \in \mathbb{R}^{C \times \frac{C}{r}}$, và $\sigma(\cdot)$ là hàm Sigmoid. Bản đồ đặc trưng đầu ra $\tilde{\mathbf{X}}_c = s_c \cdot \mathbf{X}_c$ giúp mô hình khuếch đại các kênh nhạy cảm với mật độ phương tiện và kìm hãm phản hồi từ cảnh nền tĩnh.

\subsection{Cấu trúc đầu hồi quy nén đặc trưng (Bottleneck MLP)}
\label{sec:bottleneck_mlp}

\textbf{Tầng gộp trung bình toàn cục (Global Average Pooling -- GAP):}
Tại tầng cuối cùng của mạng xương sống, tensor đặc trưng $\mathbf{F} \in \mathbb{R}^{C \times H_f \times W_f}$ được tổng hợp thành vector biểu diễn $\mathbf{f} \in \mathbb{R}^C$ thông qua phép lấy trung bình không gian:
\begin{equation}
    f_c = \frac{1}{H_f \cdot W_f} \sum_{i=1}^{H_f} \sum_{j=1}^{W_f} \mathbf{F}_c(i, j) \quad (\forall c = 1, \dots, C)
\end{equation}
Phép biến đổi này không sinh thêm tham số học được, loại bỏ nguy cơ quá khớp của các tầng kết nối dày đặc cổ điển, đồng thời bảo đảm tính bất biến đối với vị trí xuất hiện của xe trên ảnh.

\textbf{Cấu trúc đầu hồi quy Bottleneck MLP:}
Từ vector ngữ nghĩa $\mathbf{f} \in \mathbb{R}^{2048}$ xuất ra từ EfficientNet-B5, đầu hồi quy được thiết kế theo cấu trúc nén thắt cổ chai:
\begin{align}
    \mathbf{h}_1 &= \text{LayerNorm}(\mathbf{W}_1 \mathbf{f} + \mathbf{b}_1), \quad &\mathbf{W}_1 \in \mathbb{R}^{128 \times 2048}, \ \mathbf{b}_1 \in \mathbb{R}^{128} \\
    \mathbf{h}_2 &= \text{ReLU}(\mathbf{h}_1) \in \mathbb{R}^{128} \\
    \hat{\mathbf{c}} &= \mathbf{W}_2 \mathbf{h}_2 + \mathbf{b}_2, \quad &\mathbf{W}_2 \in \mathbb{R}^{2 \times 128}, \ \mathbf{b}_2 \in \mathbb{R}^2
\end{align}
Tầng thắt cổ chai 128 chiều đóng vai trò như một bộ lọc thông tin (Information Bottleneck), ép mạng phải tổng hợp các đặc trưng trừu tượng phản ánh mật độ và loại bỏ các đặc trưng cục bộ thừa thãi. Tầng Layer Normalization giúp ổn định phân phối kích hoạt trước các biến thiên điều kiện ánh sáng mạnh giữa ban ngày và ban đêm.

\subsection{Hàm mất mát bền vững Huber Loss}
\label{sec:huber_loss}
Trong điều kiện quan trắc camera thực địa, dữ liệu hình ảnh chịu ảnh hưởng từ các yếu tố môi trường cực đoan (lóa đèn xe rọi thẳng vào ống kính ban đêm, giọt nước mưa làm nhòe kính hoặc camera bị rung lắc do gió lớn), dẫn đến sự xuất hiện của các điểm ngoại lai (outliers) có sai số đếm bất thường.

Nếu sử dụng hàm mất mát trung bình bình phương sai số ($L_2$ Loss):
\begin{equation}
    \mathcal{L}_2(e) = \frac{1}{2} e^2, \quad \frac{\partial \mathcal{L}_2}{\partial e} = e
\end{equation}
gradient tăng tuyến tính theo sai số $e = c - \hat{c}$. Khi gặp một điểm ngoại lai lớn ($|e| \gg 1$), gradient bùng nổ sẽ làm biến dạng các trọng số học được, gây mất ổn định mô hình. Ngược lại, hàm mất mát sai số tuyệt đối ($L_1$ Loss) có đạo hàm không liên tục tại $e=0$, gây dao động số học quanh điểm tối ưu.

Mô hình áp dụng hàm mất mát \textbf{Huber Loss} với ngưỡng chuyển tiếp $\delta = 1.0$:
\begin{equation}
    \mathcal{L}_{\mathrm{Huber}}(e) =
    \begin{cases}
    \frac{1}{2} e^2, & \text{khi } |e| < \delta, \\[4pt]
    \delta \left( |e| - \frac{1}{2} \delta \right), & \text{khi } |e| \ge \delta
    \end{cases}
\end{equation}
Hàm ảnh hưởng (Influence Function -- đạo hàm bậc nhất của hàm mất mát) có dạng:
\begin{equation}
    \psi(e) = \frac{\partial \mathcal{L}_{\mathrm{Huber}}}{\partial e} =
    \begin{cases}
    e, & \text{khi } |e| < \delta, \\[4pt]
    \delta \cdot \text{sign}(e), & \text{khi } |e| \ge \delta
    \end{cases}
\end{equation}
Ở lân cận sai số nhỏ ($|e| < \delta$), hàm cư xử như $L_2$ Loss cho phép gradient triệt tiêu mượt mà về 0, bảo đảm tính hội tụ chính xác. Ở vùng sai số lớn ($|e| \ge \delta$), gradient bị chặn trên bởi hằng số $\pm \delta$, cô lập tác động phá hủy của các điểm dị biệt.

Hàm mất mát tổng thể để tối ưu hóa Giai đoạn 1 là:
\begin{equation}
    \mathcal{L}_{\text{Stage 1}} = \frac{1}{2} \left( \mathcal{L}_{\mathrm{Huber}}(c_{\text{car}} - \hat{c}_{\text{car}}) + \mathcal{L}_{\mathrm{Huber}}(c_{\text{bike}} - \hat{c}_{\text{bike}}) \right)
\end{equation}

% ==============================================================================
\section{Mạng nơ-ron đồ thị không -- thời gian nhận biết thời gian TA-STGCN}
\label{sec:kien_truc_ta_stgcn}

\subsection{Tổng quan bài toán dự báo dòng giao thông không -- thời gian trên đồ thị}
\label{sec:bieu_dien_khong_gian_thoi_gian}

Dòng giao thông trên mạng lưới đường bộ đô thị vận hành trong không gian phi Euclid có cấu trúc liên kết hình học phức tạp. Sự tương tác giữa các luồng xe không chỉ phụ thuộc vào khoảng cách địa lý Euclid thuần túy mà bị chi phối mạnh mẽ bởi cấu trúc hình học của đường (chiều lưu thông, dải phân cách, số làn và các nút giao cắt). 

Mạng lưới giao thông gồm $N$ vị trí quan sát được biểu diễn dưới dạng một đồ thị có hướng, có trọng số $\mathcal{G} = (\mathcal{V}, \mathcal{E}, \mathbf{W})$, trong đó:
\begin{itemize}
    \item $\mathcal{V} = \{v_1, v_2, \dots, v_N\}$ là tập hợp $N$ đỉnh đồ thị, tương ứng với $N$ vị trí cảm biến / camera thu thập dữ liệu trên mạng lưới.
    \item $\mathcal{E} \subseteq \mathcal{V} \times \mathcal{V}$ là tập hợp các cạnh có hướng, thể hiện hành lang lưu thông trực tiếp từ nút $v_i$ sang nút $v_j$.
    \item $\mathbf{W} \in \mathbb{R}^{N \times N}$ là ma trận trọng số không gian. Phần tử $W_{ij} > 0$ biểu thị mức độ liên kết không gian giữa hai nút $v_i$ và $v_j$ dựa trên khoảng cách mạng lưới thực tế; $W_{ij} = 0$ nếu không tồn tại kết nối trực tiếp.
\end{itemize}

Tại mỗi bước thời gian $t$, trạng thái giao thông trên toàn mạng lưới được mô tả bởi ma trận $\mathbf{X}_t \in \mathbb{R}^{N \times C}$, trong đó $C$ là số lượng thuộc tính quan trắc (chẳng hạn lưu lượng đếm được của từng loại xe). Xét cửa sổ quan sát lịch sử gồm $T$ bước thời gian liên tiếp, dữ liệu đầu vào của bài toán là tensor ba chiều:
\begin{equation}
    \mathcal{X} = [\mathbf{X}_{t-T+1}, \mathbf{X}_{t-T+2}, \dots, \mathbf{X}_t] \in \mathbb{R}^{T \times N \times C}
\end{equation}

Mục tiêu của bài toán dự báo không -- thời gian là tìm ánh xạ hàm học được $\mathcal{F}_\Theta$ tham số hóa bởi $\Theta$:
\begin{equation}
    \hat{\mathcal{Y}} = [\hat{\mathbf{X}}_{t+1}, \hat{\mathbf{X}}_{t+2}, \dots, \hat{\mathbf{X}}_{t+H}] = \mathcal{F}_\Theta(\mathcal{X}; \mathcal{G}) \in \mathbb{R}^{H \times N \times C}
\end{equation}
nhằm dự báo chính xác trạng thái giao thông tại toàn bộ $N$ nút trên mạng lưới qua $H$ bước thời gian tiếp theo trong tương lai.

\subsection{Nền tảng giải tích đồ thị phổ và toán tử tích chập Chebyshev}
\label{sec:ly_thuyet_do_thi_pho}

Để trích xuất các tương quan không gian phi Euclid trên đồ thị $\mathcal{G}$, lý thuyết đồ thị phổ (Spectral Graph Theory) cung cấp cơ sở giải tích cho phép định nghĩa phép biến đổi Fourier và tích chập trên miền phổ của toán tử vi phân đồ thị.

\textbf{1. Toán tử Laplacian đối xứng chuẩn hóa:}
Cho ma trận trọng số $\mathbf{W} \in \mathbb{R}^{N \times N}$ và ma trận bậc đường chéo $\mathbf{D} \in \mathbb{R}^{N \times N}$ với $D_{ii} = \sum_{j=1}^N W_{ij}$. Toán tử Laplacian đối xứng chuẩn hóa $\mathbf{L} \in \mathbb{R}^{N \times N}$ được định nghĩa:
\begin{equation}
    \mathbf{L} = \mathbf{I}_N - \mathbf{D}^{-1/2} \mathbf{W} \mathbf{D}^{-1/2}
\end{equation}
trong đó $\mathbf{I}_N$ là ma trận đơn vị cấp $N$.

\textbf{Các đặc tính đại số quan trọng của $\mathbf{L}$:}
\begin{enumerate}
    \item $\mathbf{L}$ là ma trận thực, đối xứng và nửa xác định dương. Với mọi vector tín hiệu đồ thị $\mathbf{x} \in \mathbb{R}^N$:
    \begin{equation}
        \mathbf{x}^T \mathbf{L} \mathbf{x} = \frac{1}{2} \sum_{i=1}^N \sum_{j=1}^N W_{ij} \left( \frac{x_i}{\sqrt{D_{ii}}} - \frac{x_j}{\sqrt{D_{jj}}} \right)^2 \ge 0
    \end{equation}
    Do đó, toàn bộ các giá trị riêng của $\mathbf{L}$ đều là số thực không âm.
    \item Phổ giá trị riêng của $\mathbf{L}$ bị chặn trong đoạn hữu hạn:
    \begin{equation}
        0 = \lambda_1 \le \lambda_2 \le \dots \le \lambda_N \le 2
    \end{equation}
\end{enumerate}

Nhờ tính đối xứng thực, ma trận $\mathbf{L}$ có thể phân rã phổ trực chuẩn:
\begin{equation}
    \mathbf{L} = \mathbf{U} \boldsymbol{\Lambda} \mathbf{U}^T
\end{equation}
với $\mathbf{U} = [\mathbf{u}_1, \mathbf{u}_2, \dots, \mathbf{u}_N] \in \mathbb{R}^{N \times N}$ là ma trận trực giao chứa các vector riêng thỏa mãn $\mathbf{U} \mathbf{U}^T = \mathbf{I}_N$, và $\boldsymbol{\Lambda} = \text{diag}(\lambda_1, \dots, \lambda_N)$ là ma trận đường chéo chứa các giá trị riêng tương ứng.

\textbf{2. Biến đổi Fourier và định lý tích chập trên đồ thị:}
Các vector riêng $\mathbf{u}_k$ đóng vai trò như hệ hàm sóng cơ sở trên đồ thị. Phép biến đổi Fourier trên đồ thị của tín hiệu $\mathbf{x} \in \mathbb{R}^N$ và biến đổi nghịch đảo được xác định bởi:
\begin{equation}
    \hat{\mathbf{x}} = \mathbf{U}^T \mathbf{x}, \quad \mathbf{x} = \mathbf{U} \hat{\mathbf{x}}
\end{equation}
Theo định lý tích chập, phép tích chập giữa tín hiệu $\mathbf{x}$ và bộ lọc tham số $g_\theta$ trên đồ thị là phép nhân từng phần tử trong miền tần số:
\begin{equation}
    \mathbf{x} \ast_{\mathcal{G}} g_\theta = \mathbf{U} \left( (\mathbf{U}^T \mathbf{x}) \odot (\mathbf{U}^T g_\theta) \right) = \mathbf{U} g_\theta(\boldsymbol{\Lambda}) \mathbf{U}^T \mathbf{x}
\end{equation}
trong đó $g_\theta(\boldsymbol{\Lambda}) = \text{diag}(g_\theta(\lambda_1), \dots, g_\theta(\lambda_N))$ là ma trận đáp ứng phổ của bộ lọc.

\textit{Hạn chế tính toán:} Việc tính toán trực tiếp công thức trên đòi hỏi phép phân rã phổ của $\mathbf{L}$ với độ phức tạp $\mathcal{O}(N^3)$ và phép nhân ma trận dày đặc $\mathbf{U}$ với độ phức tạp $\mathcal{O}(N^2)$, hoàn toàn không khả thi trên các mạng lưới đô thị có số nút lớn.

\textbf{3. Xấp xỉ đa thức Chebyshev cắt cụt (ChebNet):}
Để triệt tiêu nút thắt tính toán trên, Defferrard et al. \cite{chebnet} đề xuất xấp xỉ bộ lọc phổ $g_\theta(\boldsymbol{\Lambda})$ bằng chuỗi đa thức trực giao Chebyshev loại I cắt cụt bậc $K-1$.

Đa thức Chebyshev $T_k(x)$ được định nghĩa trên đoạn $[-1, 1]$ thông qua công thức truy hồi:
\begin{equation}
    T_0(x) = 1, \quad T_1(x) = x, \quad T_k(x) = 2x T_{k-1}(x) - T_{k-2}(x) \quad (\forall k \ge 2)
\end{equation}
Thực hiện phép biến đổi affine đưa phổ trị riêng $[0, \lambda_{\max}]$ về $[-1, 1]$:
\begin{equation}
    \tilde{\mathbf{L}} = \frac{2}{\lambda_{\max}} \mathbf{L} - \mathbf{I}_N
\end{equation}
Bộ lọc phổ được xấp xỉ qua $K$ số hạng đầu tiên:
\begin{equation}
    g_\theta \ast_{\mathcal{G}} \mathbf{x} \approx \sum_{k=0}^{K-1} \theta_k T_k(\tilde{\mathbf{L}}) \mathbf{x}
\end{equation}
với $\theta \in \mathbb{R}^K$ là vector các tham số học được của bộ lọc.

\textbf{Định lý 2.1 (Tính định vị cục bộ $(K-1)$-hop):} \textit{Toán tử tích chập ChebNet bậc $K-1$ đảm bảo tính định vị không gian trong phạm vi đúng $K-1$ bước nhảy lân cận trên đồ thị.}
\begin{proof}
Ma trận $\tilde{\mathbf{L}}$ là hàm bậc nhất của ma trận kề $\mathbf{W}$. Với hai nút $i$ và $j$ bất kỳ, nếu khoảng cách ngắn nhất trên đồ thị $d_{\mathcal{G}}(i, j) > k$, thì phần tử $[\tilde{\mathbf{L}}^k]_{ij} = 0$. Vì $T_k(\tilde{\mathbf{L}})$ là đa thức đại số bậc $k$ của $\tilde{\mathbf{L}}$, nó được biểu diễn thành tổ hợp tuyến tính của các lũy thừa $\tilde{\mathbf{L}}^m$ ($m \le k$). Do đó, $[T_k(\tilde{\mathbf{L}})]_{ij} = 0$ khi $d_{\mathcal{G}}(i, j) > k$. Với bậc xấp xỉ $K=3$, thông tin giao thông chỉ lan truyền trong phạm vi 2-hop lân cận, hoàn toàn tương thích với động học lan truyền sóng giao thông thực tế.
\end{proof}

\textbf{Định lý 2.2 (Độ phức tạp thời gian tuyến tính):} \textit{Phép tích chập ChebNet có độ phức tạp thời gian $\mathcal{O}(K \cdot |\mathcal{E}|)$, tuyến tính theo số cạnh của đồ thị.}
\begin{proof}
Đại lượng $\mathbf{z}_k = T_k(\tilde{\mathbf{L}}) \mathbf{x}$ được tính toán đệ quy: $\mathbf{z}_0 = \mathbf{x}$, $\mathbf{z}_1 = \tilde{\mathbf{L}}\mathbf{x}$, $\mathbf{z}_k = 2\tilde{\mathbf{L}}\mathbf{z}_{k-1} - \mathbf{z}_{k-2}$. Mỗi bước đệ quy chỉ gồm phép nhân ma trận thưa $\tilde{\mathbf{L}}$ với vector $\mathbf{z}_{k-1}$. Với đồ thị mạng lưới đường bộ có cấu trúc thưa ($|\mathcal{E}| \ll N^2$), toàn bộ quá trình tính toán $K$ bước chỉ tiêu tốn chi phí $\mathcal{O}(K \cdot |\mathcal{E}|)$, loại bỏ hoàn toàn chi phí $\mathcal{O}(N^3)$ của phép phân rã phổ.
\end{proof}

\subsection{Tích chập thời gian có cổng (Temporal Convolution with GLU)}
\label{sec:tich_chap_thoi_gian_co_cong}

Để mô hình hóa chiều thời gian trong dữ liệu chuỗi giao thông, các mô hình hồi quy tuần hoàn (RNN, LSTM, GRU) thường gặp hai hạn chế kỹ thuật:
\begin{enumerate}
    \item \textbf{Nút thắt tính toán tuần tự (Sequential Bottleneck):} Việc cập nhật trạng thái ẩn $\mathbf{h}_t = f(\mathbf{h}_{t-1}, \mathbf{x}_t)$ buộc phải thực hiện tuần tự qua từng bước thời gian, không thể khai thác kiến trúc xử lý song song khối lượng lớn của GPU.
    \item \textbf{Suy thoái gradient (Gradient Degradation):} Việc nhân chuỗi liên tiếp các ma trận trọng số qua nhiều bước thời gian dài khiến gradient dễ bị triệt tiêu hoặc bùng nổ trong quá trình lan truyền ngược qua thời gian (BPTT).
\end{enumerate}

Để khắc phục nhược điểm này, mô hình áp dụng phép tích chập thời gian 1D kết hợp cơ chế \textbf{Cổng tuyến tính (Gated Linear Units -- GLU)}.

Xét tensor tín hiệu của một nút dọc theo trục thời gian $\mathbf{X} \in \mathbb{R}^{T \times C_{\text{in}}}$. Tầng tích chập thời gian áp dụng hai bộ lọc tích chập 1D độc lập $\boldsymbol{\Theta}_1, \boldsymbol{\Theta}_2 \in \mathbb{R}^{K_t \times C_{\text{in}} \times C_{\text{out}}}$ với độ dài kernel thời gian $K_t$:
\begin{equation}
    \Gamma \ast_{\mathcal{T}} \mathbf{X} = (\mathbf{X} \star \boldsymbol{\Theta}_1 + \mathbf{b}_1) \odot \sigma(\mathbf{X} \star \boldsymbol{\Theta}_2 + \mathbf{b}_2)
\end{equation}
trong đó $\star$ là phép tích chập 1D dọc trục thời gian; $\odot$ là phép nhân Hadamard từng phần tử; và $\sigma(z) = (1 + e^{-z})^{-1}$ là hàm Sigmoid đóng vai trò van cổng đóng mở dòng thông tin.

\textbf{Phân tích giải tích luồng gradient của GLU:}
Đặt $\mathbf{P} = \mathbf{X} \star \boldsymbol{\Theta}_1 + \mathbf{b}_1$ là nhánh tuyến tính và $\mathbf{Q} = \mathbf{X} \star \boldsymbol{\Theta}_2 + \mathbf{b}_2$ là nhánh điều khiển cổng. Đầu ra của cổng là $\mathbf{Y} = \mathbf{P} \odot \sigma(\mathbf{Q})$. Đạo hàm của $\mathbf{Y}$ theo nhánh tuyến tính $\mathbf{P}$ là:
\begin{equation}
    \frac{\partial \mathbf{Y}}{\partial \mathbf{P}} = \sigma(\mathbf{Q})
\end{equation}
Khi van cổng mở ($\sigma(\mathbf{Q}) \to 1$), gradient $\frac{\partial \mathcal{L}}{\partial \mathbf{Y}}$ được truyền thẳng trực tiếp về $\mathbf{P}$ như một đường truyền tuyến tính, hoàn toàn không bị suy giảm bởi đạo hàm của các hàm phi tuyến tính như trong RNN thông thường. Cơ chế này cho phép huấn luyện các mạng tích chập sâu một cách ổn định, đồng thời toàn bộ chuỗi thời gian được xử lý song song với tốc độ cao.

\subsection{Cơ chế Self-Attention theo thời gian đa đầu (Multi-Head Temporal Self-Attention)}

Mặc dù tích chập thời gian GLU tính toán nhanh và ổn định, kiến trúc STGCN tiêu chuẩn chỉ sử dụng kernel thời gian cục bộ nhỏ ($K_t = 3$). Thiết kế này bộc lộ một \textbf{điểm yếu chí tử trong hệ thống Two-Stage Pipeline}:
\begin{itemize}
    \item Sai số ước lượng từ Giai đoạn 1 (do rung lắc camera, chói sáng hoặc xe buýt che khuất tức thời) tạo ra các điểm đột biến dị thường (outliers) trong chuỗi đếm.
    \item Do trường tiếp nhận thời gian cục bộ hẹp, kernel tích chập GLU không thể nhận biết được xung nhiễu này là dị biệt so với xu thế toàn cục. Xung lỗi lập tức truyền thẳng vào tầng tích chập không gian ChebNet và lan truyền sang toàn bộ các nút lân cận trên đồ thị mạng lưới đường bộ (hiện tượng \textit{Error Propagation}).
\end{itemize}

Để giải quyết vấn đề này, đề tài đề xuất tích hợp \textbf{Cơ chế Self-Attention theo thời gian đa đầu (Multi-Head Temporal Self-Attention -- MHTA)} vào khối xử lý không -- thời gian.

\textbf{1. Thiết kế chú ý thời gian tách rời (Decoupled Temporal Attention):}
Thay vì tính toán cơ chế chú ý liên hợp trên toàn bộ không gian và thời gian $\mathbb{R}^{N \times T}$ (gây bùng nổ bộ nhớ với đồ thị quy mô lớn $N = 608$), cơ chế chú ý thời gian được thực thi độc lập cho từng đỉnh $i \in \{1, \dots, N\}$ dọc theo chuỗi quan sát $T$ bước:
\begin{equation}
    \mathbf{Z}_i = \mathbf{X}_{i, :, :} \in \mathbb{R}^{T \times C}
\end{equation}
Để mô hình nhận biết được thứ tự thời gian của các bước quan sát, vector mã hóa vị trí (Positional Encoding) được cộng trực tiếp vào chuỗi đặc trưng:
\begin{equation}
    \tilde{\mathbf{Z}}_i = \mathbf{Z}_i + \mathbf{PE}
\end{equation}
với các thành phần hàm điều hòa:
\begin{equation}
    \mathbf{PE}_{(pos, 2j)} = \sin\left(\frac{pos}{10000^{2j/C}}\right), \quad \mathbf{PE}_{(pos, 2j+1)} = \cos\left(\frac{pos}{10000^{2j/C}}\right)
\end{equation}

Chuỗi đặc trưng $\tilde{\mathbf{Z}}_i$ được ánh xạ tuyến tính thành ba ma trận Truy vấn (Query), Khóa (Key) và Giá trị (Value):
\begin{equation}
    \mathbf{Q}_i = \tilde{\mathbf{Z}}_i \mathbf{W}^Q, \quad \mathbf{K}_i = \tilde{\mathbf{Z}}_i \mathbf{W}^K, \quad \mathbf{V}_i = \tilde{\mathbf{Z}}_i \mathbf{W}^V
\end{equation}
trong đó $\mathbf{W}^Q, \mathbf{W}^K \in \mathbb{R}^{C \times d_k}$ và $\mathbf{W}^V \in \mathbb{R}^{C \times d_v}$ là các ma trận trọng số học được ($d_k = d_v = C / h$).

Ma trận trọng số chú ý thời gian $\mathbf{A}_i \in \mathbb{R}^{T \times T}$ được tính bằng tích vô hướng chia tỷ lệ:
\begin{equation}
    \mathbf{A}_i = \text{Softmax}\left( \frac{\mathbf{Q}_i \mathbf{K}_i^T}{\sqrt{d_k}} \right)
\end{equation}
Đầu ra của cơ chế chú ý tại đỉnh $i$ là:
\begin{equation}
    \text{Attention}(\mathbf{Q}_i, \mathbf{K}_i, \mathbf{V}_i) = \mathbf{A}_i \mathbf{V}_i
\end{equation}

\textbf{2. Cơ chế đa đầu (Multi-Head Attention):}
Để đồng thời trích xuất nhiều dạng quan hệ phụ thuộc thời gian (chu kỳ tín hiệu đèn, xu thế lưu lượng tăng/giảm và mức nền ổn định), cơ chế chú ý được triển khai song song trên $h = 4$ đầu chú ý độc lập:
\begin{align}
    \text{head}_m^{(i)} &= \text{Attention}(\tilde{\mathbf{Z}}_i \mathbf{W}_m^Q, \tilde{\mathbf{Z}}_i \mathbf{W}_m^K, \tilde{\mathbf{Z}}_i \mathbf{W}_m^V) \quad (m = 1, \dots, h) \\
    \text{MHTA}(\mathbf{Z}_i) &= \text{concat}\left( \text{head}_1^{(i)}, \text{head}_2^{(i)}, \dots, \text{head}_h^{(i)} \right) \mathbf{W}^O
\end{align}
với $\mathbf{W}^O \in \mathbb{R}^{C \times C}$ là ma trận chiếu đầu ra.

\textbf{Ý nghĩa giải tích của MHTA:} Cơ chế Self-Attention theo thời gian đóng vai trò như một \textbf{bộ lọc thông thấp thích ứng dữ liệu (Data-driven Adaptive Low-pass Filter)}. Khi một bước thời gian xuất hiện xung lỗi đột biến (outlier từ Giai đoạn 1), vector truy vấn và khóa của bước thời gian đó sẽ có độ tương đồng thấp với toàn bộ bối cảnh lịch sử còn lại trong cửa sổ $T$. Do đó, hàm Softmax sẽ tự động phân bổ trọng số chú ý cực nhỏ cho bước lỗi này, làm mượt xung nhiễu trước khi tín hiệu được truyền vào đồ thị không gian.

\subsection{Kiến trúc tích hợp TA-ST-Conv và phân tích độ phức tạp tiệm cận}
\label{sec:kien_truc_stgcn_tieu_chuan}

\textbf{1. Cấu trúc khối TA-ST-Conv hoàn chỉnh:}
Mỗi khối xử lý không -- thời gian trong TA-STGCN được cấu hình theo quy trình tích hợp ba tầng chặt chẽ:
\begin{align}
    \mathbf{Z}^{(l)} &= \text{LayerNorm}\left( \mathbf{X}^{(l)} + \text{MHTA}(\mathbf{X}^{(l)}) \right) \\
    \mathbf{S}^{(l)} &= \text{SpatialChebConv}\left( \mathbf{Z}^{(l)}; \tilde{\mathbf{L}}, K=3 \right) \\
    \mathbf{X}^{(l+1)} &= \text{LayerNorm}\left( \mathbf{S}^{(l)} + \text{TemporalGLU}(\mathbf{S}^{(l)}) \right)
\end{align}
Bằng cách bố trí khối MHTA ngay tại cửa ngõ đầu vào của khối trước toán tử ChebNet, mọi xung nhiễu cục bộ tại từng nút đều được làm sạch trước khi diễn ra quá trình khuếch tán thông tin trên đồ thị không gian. Tầng tích chập thời gian GLU đặt ở cuối đảm nhiệm vai trò cập nhật động học và điều tiết luồng gradient.

\textbf{2. Khối xuất kết quả (Output Block):}
Sau khi tín hiệu truyền qua hai khối TA-ST-Conv liên tiếp, kích thước chiều thời gian được nén về 1 thông qua một tầng tích chập thời gian 1D có kích thước kernel bằng đúng chiều dài chuỗi còn lại. Tầng tích chập $1 \times 1$ cuối cùng ánh xạ số kênh về đúng tầm nhìn dự báo $H = 6$ bước:
\begin{equation}
    \hat{\mathcal{Y}} = \text{Conv}_{1 \times 1}\left( \text{ReLU}\left( \text{Conv}_{\text{time}}(\mathbf{X}^{(2)}) \right) \right) \in \mathbb{R}^{H \times N \times 1}
\end{equation}

\textbf{3. Đối sánh kiến trúc và phân tích độ phức tạp tiệm cận:}
Bảng \ref{tab:so_sanh_kien_truc_stgcn} tổng hợp đặc tính cấu trúc và độ phức tạp tính toán giữa TA-STGCN và các mô hình STGNN kinh điển trong y văn.

\begin{table}[htbp]
\centering
\caption{So sánh đặc tính kiến trúc và độ phức tạp tính toán giữa các mô hình STGNN.}
\label{tab:so_sanh_kien_truc_stgcn}
\footnotesize
\setlength{\tabcolsep}{4pt}
\renewcommand{\arraystretch}{1.2}
\begin{tabularx}{\textwidth}{l >{\raggedright\arraybackslash}p{2.3cm} >{\raggedright\arraybackslash}p{2.2cm} >{\raggedright\arraybackslash}p{2.8cm} >{\raggedright\arraybackslash}X}
\toprule
\textbf{Mô hình} & \textbf{Không gian} & \textbf{Thời gian} & \textbf{Độ phức tạp} & \textbf{Khả năng lọc nhiễu cảm biến} \\
\midrule
DCRNN \cite{dcrnn} & Lan truyền ngẫu nhiên & GRU tuần tự & $\mathcal{O}(K |\mathcal{E}| T + N^2)$ & Thấp (nhiễu tích lũy qua các bước hồi quy BPTT). \\
STGCN \cite{stgcn} & ChebNet ($K=3$) & Gated CNN 1D & $\mathcal{O}(K |\mathcal{E}| T)$ & Trung bình (kernel cục bộ $K_t=3$, dễ bị sốc ngoại lai). \\
ASTGCN \cite{astgcn} & ChebNet + Spatial Attn & Gated CNN + Temp Attn & $\mathcal{O}(T N^2 C + K |\mathcal{E}| T)$ & Trung bình (chú ý không gian bị nhiễu cục bộ đánh lừa). \\
Graph WaveNet \cite{wavenet} & GCN tự học & Dilated Causal Conv & $\mathcal{O}(N^2 + K |\mathcal{E}| T)$ & Trung bình (không có cơ chế tự làm mượt trên trục thời gian). \\
\textbf{TA-STGCN (Đề xuất)} & \textbf{ChebNet ($K=3$)} & \textbf{MHTA + GLU} & $\mathcal{O}(N T^2 C + K |\mathcal{E}| T)$ & \textbf{Rất cao (cơ chế chú ý hoạt động như bộ lọc thông thấp thích ứng dữ liệu).} \\
\bottomrule
\end{tabularx}
\end{table}

\textbf{Ưu thế vượt trội về độ phức tạp của cơ chế chú ý tách rời:}
\begin{itemize}
    \item \textbf{Cơ chế chú ý liên hợp không -- thời gian (Joint Attention):} Nếu áp dụng cơ chế chú ý trên toàn bộ không gian tensor $(N \times T)$, ma trận chú ý có kích thước $(N \cdot T) \times (N \cdot T)$ với độ phức tạp $\mathcal{O}((N \cdot T)^2 \cdot C)$. Với $N = 608, T = 24$, số phần tử ma trận lên tới $(14.592)^2 \approx 2.13 \times 10^8$, gây tràn bộ nhớ GPU ngay từ bước lan truyền xuôi đầu tiên.
    \item \textbf{Cơ chế chú ý không gian (Spatial Attention):} Việc tính toán tương quan toàn thể giữa mọi cặp nút tại mỗi bước thời gian khiến độ phức tạp tăng bậc hai theo số lượng đỉnh $\mathcal{O}(T \cdot N^2 \cdot C)$, tạo thành điểm nghẽn nghiêm trọng khi mở rộng quy mô mạng lưới.
    \item \textbf{Cơ chế chú ý thời gian tách rời trong TA-STGCN:} Do phép chú ý được thực thi độc lập và song song trên từng đỉnh dọc theo trục thời gian, độ phức tạp của thành phần chú ý chỉ là $\mathcal{O}(N \cdot T^2 \cdot C)$, tăng tuyến tính theo số đỉnh đồ thị $N$. Trong các mạng lưới giao thông thực tế, số lượng nút quan trắc luôn lớn hơn rất nhiều so với độ dài cửa sổ thời gian lịch sử ($N \gg T$, cụ thể $608 \gg 24$). Do đó, thiết kế này bảo đảm khả năng mở rộng quy mô vượt trội và tiết kiệm bộ nhớ phần cứng tối đa.
\end{itemize}
